\documentclass[11pt,a4paper]{article}
\usepackage[utf8]{inputenc}
\usepackage[T1]{fontenc}
\usepackage{lmodern}
\usepackage[french]{babel}
\usepackage{amsmath,amssymb,mathtools}
\usepackage{icomma}
\usepackage{xcolor}
\usepackage{tikz}
\usepackage{pgfplots}
\usepackage[most]{tcolorbox}
\usepackage{enumitem}
\usepackage{fancyhdr}
\usepackage[hidelinks]{hyperref}
\usepackage{titlesec}
\usepackage[a4paper,margin=2.2cm,headheight=26pt,headsep=10pt]{geometry}
\usetikzlibrary{arrows.meta,positioning,calc,intersections,angles,quotes}
\usepgfplotslibrary{fillbetween}
\pgfplotsset{compat=1.18}
\definecolor{primary}{RGB}{30,75,145}
\definecolor{primarydark}{RGB}{20,50,100}
\definecolor{defcol}{RGB}{0,114,178}
\definecolor{thmcol}{RGB}{0,135,110}
\definecolor{excol}{RGB}{0,130,85}
\definecolor{remarkcol}{RGB}{120,70,180}
\definecolor{attncol}{RGB}{200,40,55}
\definecolor{methcol}{RGB}{85,85,100}
\definecolor{areacol}{RGB}{120,160,230}
\definecolor{areacol2}{RGB}{230,150,80}
\newcommand{\dd}{\mathrm{d}}
\newcommand{\R}{\mathbb{R}}
\newcommand{\ua}{\,\text{u.a.}}
\newcommand{\tg}{\tan}\newcommand{\cotg}{\cot}
\tcbset{boxbase/.style={enhanced,breakable,boxrule=0.9pt,arc=2.5pt,left=3mm,right=3mm,top=2.3mm,bottom=2.3mm,fonttitle=\bfseries,coltitle=white,toptitle=0.6mm,bottomtitle=0.6mm}}
\newtcolorbox{methode}[1][]{boxbase,colback=methcol!7,colframe=methcol,sharp corners,title={\textsc{Comment faire}\ifx\relax#1\relax\relax\else\textmd{ : \itshape #1}\fi}}
\newtcolorbox{exemple}[1][]{boxbase,colback=excol!5,colframe=excol,title={\textsc{Exemple}\ifx\relax#1\relax\relax\else\textmd{ : \itshape #1}\fi}}
\newtcolorbox{idee}[1][]{boxbase,colback=defcol!5,colframe=defcol,title={\textsc{Idée clé}\ifx\relax#1\relax\relax\else\textmd{ : \itshape #1}\fi}}
\newtcolorbox{remarque}[1][]{boxbase,colback=remarkcol!6,colframe=remarkcol,title={\textsc{Remarque}\ifx\relax#1\relax\relax\else\textmd{ : \itshape #1}\fi}}
\newtcolorbox{attention}[1][]{boxbase,colback=attncol!6,colframe=attncol,title={\textsc{Attention}\ifx\relax#1\relax\relax\else\textmd{ : \itshape #1}\fi}}
\newtcolorbox{exo}[1][]{boxbase,colback={primary!4},colframe=primary,title={\textsc{Exercice #1}}}
\newtcolorbox{corr}[1][]{boxbase,colback={thmcol!5},colframe=thmcol,title={\textsc{Corrigé #1}}}
\tikzset{axe/.style={->,>=Stealth,black!65,line width=0.7pt},courbe/.style={primary,line width=1.3pt},courbe2/.style={areacol2!90!black,line width=1.3pt},dvert/.style={gray,dashed,line width=0.7pt}}
\pagestyle{fancy}\fancyhf{}
\fancyhead[L]{\small\textcolor{primarydark}{\textbf{Fiche 7} — Intégrale d'une fonction continue}}
\fancyhead[R]{\small\textcolor{primarydark}{\textsc{Thème 2 — Intégration par parties}}}
\fancyfoot[C]{\small\thepage}
\renewcommand{\headrulewidth}{0.8pt}
\titleformat{\section}{\large\bfseries\color{primarydark}}{\thesection}{0.6em}{}
\begin{document}
\section*{Niveau Facile — Corrigé}

\begin{corr}[1]
$\displaystyle\int x\,e^x\,\dd x$. \textbf{Choix :} $U=x$, $V'=e^x$, donc $U'=1$, $V=e^x$.
\[
\int x\,e^x\,\dd x=x\,e^x-\int e^x\cdot 1\,\dd x=x\,e^x-e^x+C=\boxed{(x-1)\,e^x+C.}
\]
\textbf{Vérif.} $\bigl[(x-1)e^x\bigr]'=e^x+(x-1)e^x=x\,e^x$. \checkmark
\end{corr}

\begin{corr}[2]
$\displaystyle\int x\cos x\,\dd x$. \textbf{Choix :} $U=x$, $V'=\cos x$, donc $U'=1$, $V=\sin x$.
\[
\int x\cos x\,\dd x=x\sin x-\int \sin x\cdot 1\,\dd x=x\sin x-(-\cos x)+C=\boxed{x\sin x+\cos x+C.}
\]
\textbf{Vérif.} $\bigl[x\sin x+\cos x\bigr]'=\sin x+x\cos x-\sin x=x\cos x$. \checkmark
\end{corr}

\begin{corr}[3]
$\displaystyle\int x\sin x\,\dd x$. \textbf{Choix :} $U=x$, $V'=\sin x$, donc $U'=1$, $V=-\cos x$.
\[
\int x\sin x\,\dd x=x(-\cos x)-\int(-\cos x)\cdot 1\,\dd x=-x\cos x+\int\cos x\,\dd x=\boxed{-x\cos x+\sin x+C.}
\]
\textbf{Vérif.} $\bigl[-x\cos x+\sin x\bigr]'=-\cos x+x\sin x+\cos x=x\sin x$. \checkmark
\end{corr}

\begin{corr}[4]
$\displaystyle\int x\,e^{-x}\,\dd x$. \textbf{Choix :} $U=x$, $V'=e^{-x}$, donc $U'=1$, $V=-e^{-x}$.
\[
\int x\,e^{-x}\,\dd x=x(-e^{-x})-\int(-e^{-x})\,\dd x=-x\,e^{-x}+\int e^{-x}\,\dd x=-x\,e^{-x}-e^{-x}+C=\boxed{-(x+1)\,e^{-x}+C.}
\]
\textbf{Vérif.} $\bigl[-(x+1)e^{-x}\bigr]'=-e^{-x}-(x+1)(-e^{-x})=-e^{-x}+(x+1)e^{-x}=x\,e^{-x}$. \checkmark
\end{corr}

\begin{corr}[5]
$\displaystyle\int (x+2)\,e^{x}\,\dd x$. \textbf{Choix :} $U=x+2$, $V'=e^x$, donc $U'=1$, $V=e^x$.
\[
\int (x+2)e^x\,\dd x=(x+2)e^x-\int e^x\,\dd x=(x+2)e^x-e^x+C=\boxed{(x+1)\,e^{x}+C.}
\]
\textbf{Vérif.} $\bigl[(x+1)e^x\bigr]'=e^x+(x+1)e^x=(x+2)e^x$. \checkmark
\end{corr}
\end{document}
